For the arm-normalized profile, strict positivity and finiteness of every \(w_{d,\mathrm{arm}}\) follow from positive costs, Assumption~\(\mathrm{ass{:}nondegenerate\text{-}components}\), and Theorem~\(\mathrm{thm{:}budget\text{-}profile}\).  For the IKM-normalized profile, Assumption~\(\mathrm{ass{:}source\text{-}arm\text{-}positivity}\) and Definition~\(\mathrm{def{:}ikmw\text{-}normalized\text{-}design\text{-}value}\) give
\[
 V_{E,d}^{\mathrm{IKM}}=(1-p_{E,1})^2V_{E,d}>0,
 \qquad
 V_{O,d}^{\mathrm{IKM}}=(1-p_{O,1})^2V_{O,d}>0.
\]
They are finite by Assumption~\(\mathrm{ass{:}nondegenerate\text{-}components}\); positive finite costs then make every \(w_{d,\mathrm{IKM}}\) finite and strictly positive.

Fix \((j,s)\).  The finite set \(\mathcal O_{j,P}\) contains every diagonal pair, so its maximum exists and is at least one.  If the pairwise implication holds with \(C\), division by the positive \(w_{d',s}\) and maximization show \(C_{j,s}^{\mathrm{ord}}(P)\le C\).  Conversely, the definition of the maximum shows that \(C\ge C_{j,s}^{\mathrm{ord}}(P)\) implies every pairwise inequality.  Thus \(C_{j,s}^{\mathrm{ord}}(P)\) is sharp.

Choose \(d_j\in\arg\min_{d\in\mathcal D_j(P)}w_{d,s}\) and \(d_s\in\arg\min_{d\in\mathcal D}w_{d,s}\).  Since \(d_j\) minimizes \(p_{d,j}\), one has \(p_{d_j,j}\le p_{d_s,j}\).  Consequently
\[
 1\le\frac{w_{d_j,s}}{w_{d_s,s}}
 =C_{j,s}^{\mathrm{sel}}(P)
 \le C_{j,s}^{\mathrm{ord}}(P).
\]
The same division argument proves that the middle ratio is the least best-tie constant.  Taking suprema proves the asserted uniform equivalences.  Nothing in this argument depends on whether \(p_{d,j}\) is the unadjusted risk or the treatment-aware risk, so it establishes both rules simultaneously.

Under the numerical-handle hypotheses, the definitions of the common infimum and supremum give
\[
 \Lambda_{j,s,-}(P)Q_{R,d}^{j}(x)
 \le F_{d,s}(x)
 \le\Lambda_{j,s,+}(P)Q_{R,d}^{j}(x)
\]
on every stated cone.  If \((d,d')\in\mathcal O_{j,P}\), evaluation at the transported directions yields
\[
\begin{aligned}
 w_{d,s}
 &\le\Lambda_{j,s,+}(P)p_{d,j}
 \le\Lambda_{j,s,+}(P)p_{d',j}\\
 &\le\frac{\Lambda_{j,s,+}(P)}{\Lambda_{j,s,-}(P)}w_{d',s}.
\end{aligned}
\]
Maximizing proves the certificate.  The infimum and supremum are respectively the largest possible common lower-envelope constant and the smallest possible common upper-envelope constant, which proves the claimed weakness.

Finally, for \(u=c_{E,d}Q_{E,d}^{s}(x)\) and \(w=c_{O,d}Q_{O,d}^{s}(x)\),
\[
 \frac{u}{\alpha}+\frac{w}{1-\alpha}-(\sqrt u+\sqrt w)^2
 =\frac{\{(1-\alpha)\sqrt u-\alpha\sqrt w\}^{2}}
 {\alpha(1-\alpha)}\ge0.
\]
Equality holds at the interior optimizer when \(u,w>0\), and endpoint limits prove the identity when a component vanishes.  Thus the profiled formula follows.  Taking \(Q_E(x)=x_1^2\), \(Q_O(x)=x_2^2\), and unit costs gives \(F(x)=(|x_1|+|x_2|)^2\), which fails the parallelogram identity.  Hence profiling is not generally quadratic.  A single largest generalized eigenvalue supplies only an upper fixed-share envelope; the displayed selection comparison also needs a positive lower envelope and saturation by raw-risk-ordered transported directions.  This proves the result.